% !PRINTED 674
\weberpage{674}

%<<CONT
$\mathfrak{A}$ äquivalenten Nenner und dasselbe Konstantensystem erzeugt,
und dies Polygon $\mathfrak{A}'$ ist also nur abhängig von den Konstanten
und von den Polygonen $\mathfrak{A}_1, \mathfrak{A}_2, \ldots, \mathfrak{A}_s$. Der Inbegriff dieser
Polygone wird eine \gesperrt{Polygonschar mit der Basis} $\mathfrak{A}_1, \mathfrak{A}_2, \ldots, \mathfrak{A}_s$
genannt und mit

\weq{9}{(\mathfrak{A}_1, \mathfrak{A}_2, \ldots, \mathfrak{A}_s)}

bezeichnet.

Wenn die Polygone $\mathfrak{A}_1, \mathfrak{A}_2, \ldots, \mathfrak{A}_s$ den größten gemeinschaftlichen
 Teiler $\mathfrak{M}$ haben, so ist dieser Teiler in allen Polygonen
der Schar enthalten und heißt der \gesperrt{Teiler der Schar}. Aber
man kann in der Schar ein Polygon $\mathfrak{A}' = \mathfrak{M}\mathfrak{B}$ so bestimmen,
daß $\mathfrak{B}$ beliebig gegebene Punkte nicht enthält.

Denn ist ein Punkt $\mathfrak{P}$ $\mu$ mal in $\mathfrak{M}$ und $\nu$ mal in $\mathfrak{A}$ enthalten
und $\varrho$ eine Funktion, die in $\mathfrak{P}$ von der ersten Ordnung verschwindet,
so ist in den Entwickelungen (2):

\[
m = \mu - \nu,
\]

und die Entwickelung von $\eta = \mathfrak{A}' : \mathfrak{A} = \mathfrak{M}\mathfrak{B} : \mathfrak{A}$ lautet

\[
\eta = e\varrho^m + \sigma\varrho^{m+1},
\]

worin

\[
e = e_1c_1 + e_2c_2 + \cdots + e_sc_s.
\]

Wählt man also die $c_i$ so, daß $e$ nicht verschwindet, so ist
$\mathfrak{P}$ nicht in $\mathfrak{B}$ enthalten.

Eine Schar, die einen Teiler $\mathfrak{M}$ hat, heißt eine \gesperrt{uneigentliche
 Schar}.

Wenn zwischen den Funktionen $\eta_1, \eta_2, \ldots, \eta_s$ eine Relation
von der Form besteht:

\weq{10}{c_1\eta_1 + c_2\eta_2 + \cdots + c_s\eta_s = 0,}

in der die Konstanten $c_1, c_2, \ldots, c_s$ nicht alle $= 0$ sind, so besteht
 dieselbe Relation

\weq{11}{c_1\eta_1' + c_2\eta_1' + \cdots + c_s\eta_s' = 0}
% [sic] printed c_2\eta_1' as the second term, should be c_2\eta_2': (11) is (10) multiplied by \xi, with \xi\eta_2 = \eta_2' by (7) and as in (8)

zwischen den Funktionen (7). Wir nennen dann diese Funktionen
\gesperrt{linear abhängig}. Dem entsprechend heißen auch die Polygone
der Schar $\mathfrak{A}_1, \mathfrak{A}_2, \ldots, \mathfrak{A}_s$ linear abhängig. Man kann dann eine
der Funktionen, etwa $\eta_s$, linear durch die übrigen ausdrücken,
und die Schar $(\eta_1, \eta_2, \ldots, \eta_s)$ ist mit $(\eta_1, \eta_2, \ldots, \eta_{s-1})$ identisch.
Gleiches gilt von den Polygonscharen $(\mathfrak{A}_1, \mathfrak{A}_2, \ldots, \mathfrak{A}_s)$ und $(\mathfrak{A}_1,
\mathfrak{A}_2, \ldots, \mathfrak{A}_{s-1})$. Durch wiederholte Anwendung dieser Reduktion
kann man jede Schar durch eine \gesperrt{linear unabhängige oder
irreduzible Basis} darstellen. Zwei irreduzible Basen einer
%CONT>>
